\DocumentMetadata{lang=en-US, testphase=phase-II}
\documentclass[11pt, letterpaper]{article}
\input{../../shared/preamble}
\input{../../shared/syllabus-format}

\begin{document}
\selectlanguage{english}
\addto\captionsenglish{\def\figurename{Figure}}

%-------------------------------
% MASTHEAD
%-------------------------------
\epmasthead{ENGINEERING PHYSICS \textbullet\ CAPSTONE}{EP Capstone}{A Concentration of the Physics Major}

\begin{infocard}
term       \textbar\ Fall 2026~--~Spring 2027 (year-long)\\
meets      \textbar\ Mondays, 1:00--1:50~p.m.\\
where      \textbar\ Small Hall 235 \textit{\ttfamily(primary; see notes below)}\\
instructor \textbar\ Ran Yang, Ph.D. \textbar\ \href{mailto:rxyan2@wm.edu}{rxyan2@wm.edu} \textbar\ Small 252\\
site       \textbar\ \url{https://ep.yangran.org}
\end{infocard}

\vspace{0.5ex}
{\footnotesize\itshape\color{labMauveDark}Some sessions, especially guest lectures or workshops, may take place in other campus locations as announced in class or on Blackboard.}

%-------------------------------
% COURSE OVERVIEW
%-------------------------------
\section{Course Overview}
This capstone is a redesigned, year-long experience that equips Engineering Physics (EP) seniors to thrive in graduate study, industry, and national labs. The course treats AI in the broadest sense --- not only the technical craft of building systems, but the business, legal, and policy currents that decide whether an idea survives contact with the world. Those currents are inseparable from modern engineering practice, and this course is built around that connection.

\begin{tcbitemize}[
    raster columns=2, raster rows=2, raster equal height,
    raster column skip=3mm, raster row skip=3mm,
    enhanced, colback=labLavender!10, colframe=labGreen,
    boxrule=0pt, leftrule=2pt, arc=0pt,
    left=3.5mm, right=3.5mm, top=2.5mm, bottom=2.5mm,
    fonttitle=\ttfamily\bfseries\small, coltitle=white
]
\tcbitem[title={01 \textbar\ AI, TECHNICAL TO STRATEGIC}]
    \small Technical foundations and systems thinking, then the ecosystem around them --- policy, ethics, intellectual property, and entrepreneurship --- with guest experts from law, business, computing, and the library.
\tcbitem[title={02 \textbar\ RESEARCH \& LIT REVIEW WITH AI}]
    \small Teams use AI tools to accelerate literature review and research synthesis, while learning to critically evaluate and properly cite every AI-assisted source.
\tcbitem[title={03 \textbar\ CAREER \& COMMUNICATION}]
    \small A sustained focus on the soft skills that get ideas funded and hired: presentations, public speaking, interview preparation, networking, and pitch craft, alongside entrepreneurship and venture-building fundamentals.
\tcbitem[title={04 \textbar\ QUANTUM ENGINEERING}]
    \small Connecting principles to applications --- from modeling and control to hardware and measurement --- building toward the public Showcase.
\end{tcbitemize}

\medskip
Throughout the year, students develop both hard skills and transferable professional skills. They deliver mid-semester and final presentations each term, build a polished portfolio and pitch, complete career development milestones, and learn how science and engineering operate beyond the university. The result is more than a project check-in --- it is a launchpad for the next step in a student's career.

%-------------------------------
% COURSE EXPECTATIONS
%-------------------------------
\section{Course Expectations}
To complete the EP Capstone successfully, students are expected to:
\begin{itemize}
    \item Attend all Monday class sessions (1:00--1:50~p.m.) and actively participate in guest lectures and discussions.
    \item Meet with their faculty team advisor at least one hour per week.
    \item Commit at least 10 hours per week to their project. Since this is a three-credit course per semester, that workload expectation reflects both in-class and out-of-class time.
    \item \textbf{Honors students:} In addition to the capstone requirements, students pursuing honors are expected to register for an extra credit hour for their thesis and complete two credit hours of honors research across both semesters.
    \item Deliver a mid-semester team presentation each semester, plus a final team presentation each semester --- in fall, alongside a written paper, during the University final exam period; in spring, at the Spring Showcase.
    \item Complete all professional development assignments, including resume, LinkedIn profile, and pitch preparation.
    \item Participate in the Spring Showcase (April 2027), which serves as the culminating public presentation of their work.
\end{itemize}

%-------------------------------
% READINGS
%-------------------------------
\section{Readings}
Guest lecturers may assign short readings in advance of their sessions. Students are expected to complete these before class to enable informed discussion.

%-------------------------------
% COURSE MATERIALS
%-------------------------------
\section{Course Materials \& Technology}
\begin{itemize}
    \item \textbf{Gradescope} (\url{https://www.gradescope.com/}): Individual assignments (e.g., pitch write-ups, professional-development deliverables) are submitted as PDFs through Gradescope.
    \item \textbf{Personal Computer:} Bring a charged laptop with necessary adapters for sessions that involve in-class work.
\end{itemize}

%-------------------------------
% GRADING
%-------------------------------
\section{Grading}
\label{sec:grading}
\begin{itemize}
    \item At the end of the fall semester, students will receive a grade of G (continuing).
    \item A final letter grade will be assigned upon successful completion of the full year in Spring 2027, based on performance across both semesters.
    \item The final grade reflects attendance and participation, team advisor evaluations, mid-semester and final presentations, the fall paper, professional development milestones, and engagement with guest sessions.
\end{itemize}

{\small\itshape\color{labMauveDark}The week-by-week meeting schedule is published as a companion document, the EP Capstone Schedule (Fall 2026).}

%-------------------------------
% ACADEMIC CALENDAR
%-------------------------------
\section{Important Academic Calendar Dates --- Fall 2026}
\begin{itemize}
    \item \textbf{Labor Day (No classes):} September 7
    \item \textbf{Add/drop deadline:} September 4
    \item \textbf{Fall Break:} October 8--11
    \item \textbf{Midterm grading period:} October 5--25
    \item \textbf{Withdrawal deadline:} October 26
    \item \textbf{Election Day (No classes):} November 3
    \item \textbf{Remote Instruction Days:} November 23--24
    \item \textbf{Thanksgiving Break:} November 25--29
    \item \textbf{Last day of classes:} December 4
    \item \textbf{University final exam period:} December 7--11 and December 14--15
    \item \textbf{Fall final presentations \& paper due:} Wednesday, December 9, 2:00--5:00~p.m.
    \item \textbf{Fall grades due:} January 4, 2027
\end{itemize}
{\ttfamily\small\color{labMauveDark}Source: W\&M University Registrar, 2026--2027 Undergraduate Academic Calendar (\url{https://www.wm.edu/offices/registrar/catalogs-calendars-exams/ugcalendar/26-27/}).}

%-------------------------------
% POLICIES
%-------------------------------
\section{Policies}
\begin{itemize}
    \item \textbf{Attendance and Participation:} Attendance is mandatory for all Monday class sessions. Notify the instructor in advance if you must miss a class due to illness or unavoidable conflict.
    \item \textbf{Team Advisor Meetings:} Teams must meet with their faculty advisor at least one hour per week; missed advising meetings should be rescheduled promptly with the advisor.
\end{itemize}

\subsection{Attendance, Absence, and Late Work Policy}
\label{sec:attendance}
\begin{itemize}
    \item \textbf{Accessibility Accommodations.} The course will comply with all accommodations approved by Student Accessibility Services (SAS; see \hyperref[sec:ssr]{Student Success Resources}), including extended time and any other approved adjustments. Students with an SAS accommodation letter should share it with the course instructor as early in the semester as possible.
    \item \textbf{Authorization Rests with the Instructor.} Only the course instructor may excuse an absence or authorize a late submission.
    \item \textbf{Advance Notice Requirement.} A student who anticipates missing a class session must notify the instructor at least one week (seven calendar days) in advance. The absence is considered excused only once the instructor has reviewed the request and granted approval in writing (e.g., by email). Absences must be substantiated with appropriate official documentation, such as a note from a treating physician or verification from the Dean of Students Office.
    \item \textbf{Emergency Exceptions.} If a genuine emergency makes one week's advance notice unreasonable to provide, the instructor may still excuse the absence, provided the student notifies the instructor as soon as circumstances reasonably allow and can show a good-faith effort to comply with this policy.
    \item \textbf{Missed Presentations or Deliverables.} If a student misses a presentation or another major assignment, the student must contact the instructor immediately. Make-up opportunities will be handled on a case-by-case basis; an unexcused absence from a presentation is not guaranteed a make-up.
\end{itemize}

%-------------------------------
% MIDTERM & FINAL PRESENTATIONS
%-------------------------------
\section{Midterm \& Final Presentations}
\begin{itemize}
    \item \textbf{Mid-Semester Presentations:} Held in class (Week 9, October 19); every team presents project progress to the instructor and peers.
    \item \textbf{Presentation Prep:} Week 14 (November 23) is a remote team-preparation day; Week 15 (November 30) covers tips for presentations ahead of finals.
    \item \textbf{Fall Final Presentations \& Paper:} Wednesday, December 9, 2:00--5:00~p.m., during the University final exam period. Teams deliver their mid-year presentation and submit the accompanying paper. Final grades for the semester are not issued until the full year is complete (see \hyperref[sec:grading]{Grading}).
\end{itemize}

%-------------------------------
% AI POLICY
%-------------------------------
\section{Policy on Artificial Intelligence}
Students may use AI tools for learning, brainstorming, and design only with explicit permission from the instructor. Even with permission, AI-generated content must be critically evaluated, and all sources --- including AI systems --- must be cited. Using AI to complete assignments without acknowledgment or without explicit permission is a violation of the Honor Code and will be treated as plagiarism.

%-------------------------------
% BEHAVIOR
%-------------------------------
\section{Behavior and Professionalism}
\begin{itemize}
    \item Be punctual and prepared.
    \item Be respectful to guest speakers and peers.
    \item Engage actively during Q\&A and discussions.
    \item Follow W\&M's Honor Code at all times.
\end{itemize}

%-------------------------------
% FERPA & COMMUNICATION
%-------------------------------
\section{Student Privacy, FERPA, \& Communication}
In compliance with the Family Educational Rights and Privacy Act (FERPA; \url{https://www.wm.edu/offices/registrar/confidentiality-privacy/ferpa-policy/}), student grades and records cannot be shared with parents or third parties without the student's written consent. Instructors will not engage directly with parents regarding student performance.

Course announcements and materials will be posted on Blackboard. For email communication with groups of students, all recipients will be BCC'ed in accordance with privacy regulations.

%-------------------------------
% STUDENT SUCCESS
%-------------------------------
\section{Student Success Resources}
\label{sec:ssr}
\begin{itemize}
    \item \textbf{Student Accessibility Services:} \href{mailto:sas@wm.edu}{sas@wm.edu} \textbar\ 757-221-2512 \textbar\ \url{https://www.wm.edu/sas}
    \item \textbf{Dean of Students Office:} \url{https://www.wm.edu/offices/deanofstudents/}
    \item \textbf{Student Success Center:}\\ \url{https://www.wm.edu/offices/sas/student-success-center/}
    \item \textbf{Counseling Center:} \url{https://www.wm.edu/offices/wellness/counselingcenter/}
\end{itemize}

%-------------------------------
% HONOR CODE
%-------------------------------
\section{Honor Code}
William \& Mary's Honor Code (\url{https://www.wm.edu/offices/deanofstudents/services/communityvalues/honorcodeandcouncils/honorcode/}) requires honesty and integrity in all academic work. Violations include plagiarism, unauthorized collaboration, use of unauthorized materials, fabrication of data, and lying. Suspected violations will be reported to the Honor Council.

%-------------------------------
% FINAL NOTE
%-------------------------------
\section{Final Note}
The EP Capstone is designed to prepare you for your next step in industry, research, or graduate study. This fall semester is only the first stage --- the course culminates in the Spring 2027 Showcase, where teams will present their projects publicly. Full engagement throughout the year is essential for successful completion.

%-------------------------------
% SYLLABUS CHANGES
%-------------------------------
\section{Syllabus Modifications}
The instructor reserves the right to modify the contents of this syllabus, including course policies, grading criteria, and the schedule, at any time during the semester. Students will be notified of any such changes via Blackboard and/or email in a timely manner.

{\ttfamily\small\color{labMauveDark}Revised: August 23, 2026.}

%-------------------------------
% REFERENCES
%-------------------------------
\section{References}
{\small
Information in this syllabus was verified against the following official William \& Mary webpages:
\begin{itemize}
    \item University Registrar --- 2026--2027 Undergraduate Academic Calendar: \url{https://www.wm.edu/offices/registrar/catalogs-calendars-exams/ugcalendar/26-27/}
    \item University Registrar --- FERPA Policy: \url{https://www.wm.edu/offices/registrar/confidentiality-privacy/ferpa-policy/}
    \item Dean of Students --- Honor Code \& Honor Councils: \url{https://www.wm.edu/offices/deanofstudents/services/communityvalues/honorcodeandcouncils/honorcode/}
    \item Student Accessibility Services: \url{https://www.wm.edu/sas}
    \item Dean of Students Office: \url{https://www.wm.edu/offices/deanofstudents/}
    \item Student Success Center: \url{https://www.wm.edu/offices/sas/student-success-center/}
    \item Counseling Center: \url{https://www.wm.edu/offices/wellness/counselingcenter/}
    \item Information Technology --- Blackboard at W\&M: \url{https://blackboard.wm.edu/}
\end{itemize}
}

\end{document}
